\section{Additional tables for the present-day analysis}\label{sec:sitables}

\begin{table*}[!ht]
\caption{Access--number decomposition by birth cohort (General Social Survey, ages 45--79). Units and intervals as in
Table~\ref{tab:gss}. Cells with fewer than 300 respondents are omitted.}\label{tab:gsscohort}
\scriptsize
\begin{tabularx}{\textwidth}{@{}P{0.15\textwidth}llrCCC@{}}
\toprule
\tabinput{generated/tab_gss_cohort.tex}
\bottomrule
\end{tabularx}
\end{table*}

\begin{table*}[!ht]
\caption{Robustness of the decomposition for schooling and income. Intervals are from 400 bootstrap resamples.}\label{tab:gssrobust}
\scriptsize
\begin{tabularx}{\textwidth}{@{}P{0.17\textwidth}P{0.16\textwidth}CCC@{}}
\toprule
\tabinput{generated/tab_gss_robust.tex}
\bottomrule
\end{tabularx}
\end{table*}

\begin{table*}[!ht]
\caption{Three episodes of selection (General Social Survey, ages 45--79, cohorts 1900--1979): pairing is the mean standardized
trait among the ever-married; conversion is the further difference between parents and the ever-married; number is as in
Table~\ref{tab:gss}. The three sum to the total. Units: standard deviations of the trait per generation; 95\% bootstrap
intervals from 1,000 resamples.}\label{tab:gssthree}
\scriptsize
\begin{tabularx}{\textwidth}{@{}P{0.17\textwidth}lCCCC@{}}
\toprule
\tabinput{generated/tab_gss_three.tex}
\bottomrule
\end{tabularx}
\end{table*}

\begin{table*}[!ht]
\caption{The pairing gate by birth cohort (General Social Survey, ages 45--79). Never married: percentage who had never
married when interviewed. Parents among the never married: percentage of those with at least one child. Children of the
never married: their share of all children reported by respondents of that sex and cohort. Pairing differential: mean
standardized log income among the ever-married (the first term of Table~\ref{tab:gssthree}), income measured at ages
45--64. 95\% bootstrap intervals from 1,000 resamples in brackets.}\label{tab:gssgate}
\scriptsize
\begin{tabularx}{\textwidth}{@{}llCCCC@{}}
\toprule
\tabinput{generated/tab_gss_gate.tex}
\bottomrule
\end{tabularx}
\end{table*}

\begin{table*}[!ht]
\caption{Said and did in the Columbia speed-dating experiment~\cite{fisman2006}. Said: mean points of 100 given to each
attribute before the event. Did: change in the probability of a yes, in percentage points, per standard deviation of the
rating, with a fixed effect for each chooser, using the chooser's own rating or the mean rating the partner received from
everyone else (panel). Agreement: intraclass correlation between two observers of the same person; observers for 0.8 follows
Eq.~\eqref{eq:sb}. W: women choosing; M: men choosing. Shared interests belong to a pair, so no agreement is defined.}\label{tab:speed}
\scriptsize
\begin{tabularx}{\textwidth}{@{}Lrrrrrrrr@{}}
\toprule
\tabinput{generated/tab_speed.tex}
\bottomrule
\end{tabularx}
\end{table*}

\begin{table*}[!ht]
\caption{Family size by marital history (General Social Survey, ages 45--79, cohorts 1900--1979). Relative family size is
the number of children divided by the mean for the respondent's sex and decade of birth; 95\% bootstrap intervals from 500
resamples in brackets.}\label{tab:gsshist}
\scriptsize
\begin{tabularx}{\textwidth}{@{}lLrrrC@{}}
\toprule
\tabinput{generated/tab_gss_hist.tex}
\bottomrule
\end{tabularx}
\end{table*}

\begin{table*}[!ht]
\caption{Spouses' schooling compared, by birth cohort of the respondent (General Social Survey, currently married
respondents aged 25 and over). 95\% bootstrap intervals from 500 resamples in brackets.}\label{tab:gsshyp}
\scriptsize
\begin{tabularx}{\textwidth}{@{}lrCCCr@{}}
\toprule
\tabinput{generated/tab_gss_hyp.tex}
\bottomrule
\end{tabularx}
\end{table*}
